\chapter{Step 1 — Identifying Hazards and Risk Sources}
\label{chap:Step 1 — Identifying Hazards and Risk Sources}

%---------------------------- SECTION ----------------------------
\section{The Most Important Step You Can Get Wrong}
\paragraph{There is a reasonable argument that hazard identification is the single most consequential step in the entire risk assessment process. Every subsequent step — evaluating risks, designing controls, monitoring outcomes — depends entirely on the quality of what is identified here. A risk that is not identified cannot be evaluated. A risk that is not evaluated cannot be controlled. And a risk that is not controlled has a habit of materialising at the worst possible moment.}
\paragraph{The uncomfortable truth is that most significant organisational failures — regulatory breaches, financial losses, reputational crises, and operational disasters — can be traced back, at least in part, to risks that were known to someone in the organisation but were never formally surfaced, or to risks that nobody had thought to look for. The 2008 financial crisis, to return to an example from Chapter~\ref{chap:A Brief History of Risk Assessment}, was not primarily a failure of risk modelling. It was, at its root, a failure of risk identification — a collective failure to ask seriously enough: \textbf{\textit{what could really go wrong here?}}}
\paragraph{This chapter is about how to address that question quite well.}

%---------------------------- SECTION ----------------------------
\section{What We Are Looking For}
\paragraph{Before exploring the techniques for identifying hazards and risk sources, it is worth being clear about what we are actually looking for — because the answer is broader than it might initially appear.}
\paragraph{In Chapter~\ref{chap:What is Risk}, we defined a hazard as a source of potential harm — a situation, condition, action, or failure that has the potential to cause a negative outcome. In the context of a formal risk assessment, we are looking for:}
\begin{itemize}
    \bitem{Events}{specific things that could happen, whether internally driven or externally triggered, that would have a negative impact on the organisation's objectives.}
    \bitem{Conditions}{ongoing states or circumstances that create exposure, even in the absence of a specific triggering event.}
    \bitem{Behaviours}{patterns of action or inaction, by individuals or groups, that create or increase risk.}
    \bitem{Dependencies}{relationships, processes, systems, or third parties whose failure or underperformance would create risk.}
    \bitem{Changes}{planned or unplanned changes in the organisation's environment, structure, or operations that introduce new risk or alter existing exposures.}
    \bitem{Absences}{the absence of something that should be present, such as a missing control, an unfilled role, a gap in policy, or an untested process.}
\end{itemize}
\paragraph{That last category — absences — is particularly important and particularly easy to overlook. Risk identification that focuses exclusively on things that could actively go wrong may miss the quieter but equally significant risks that arise from things that are simply not being done.}

%---------------------------- SECTION ----------------------------
\section{The Conditions for Good Risk Identification}
\paragraph{Before reaching for specific techniques, it is worth reflecting on the conditions that need to be in place for risk identification to work well. Even the most sophisticated methodology will produce poor results if these conditions are absent.}

\subsection{Psychological Safety}
\paragraph{People identify risks most effectively when they feel safe doing so — when they believe that raising a concern will be welcomed rather than penalised, that identifying a problem in their own area will not be used against them, and that uncomfortable questions are genuinely valued rather than merely tolerated.}
\paragraph{The absence of psychological safety is one of the most significant and most consistently underestimated barriers to effective risk identification. In organisations where raising concerns is perceived as risky — where the messenger tends to be shot, where admitting uncertainty is seen as weakness, or where challenge to established thinking is unwelcome — risks accumulate silently until they become crises.}
\paragraph{For compliance professionals, creating the conditions for psychological safety is partly a cultural challenge — one that requires leadership commitment and consistent behaviour over time. But it is also something that can be directly influenced by how risk identification sessions are designed and facilitated. We will return to this below.}

\subsection{Diverse Perspectives}
\paragraph{No single individual or team has a complete view of an organisation's risk landscape. Effective risk identification draws on diverse perspectives — people with different roles, different levels of seniority, different functional backgrounds, and different relationships to the activities being assessed.}
\paragraph{This diversity matters for several reasons. Different vantage points reveal different risks. People at the front line of operations often see risks that senior managers are not aware of. People from different functional backgrounds bring different mental models and different categories of concern. And people who are not deeply embedded in a particular process are sometimes better placed to ask the naive but important question: \textit{why do we do it this way, and what could go wrong if we do?}}

\subsection{Structured Curiosity}
\paragraph{Good risk identification requires a combination of structure — to ensure that the process is systematic and comprehensive — and genuine curiosity — the willingness to look beyond the obvious, to challenge assumptions, and to consider scenarios that have not yet occurred.}
\paragraph{The structured element ensures that the process covers all relevant areas and does not rely on the memory or intuition of a single individual. The curiosity element ensures that the structure does not become a constraint — a checklist that is worked through mechanically without genuine engagement with what might be missed.}

%---------------------------- SECTION ----------------------------
\section{Techniques for Identifying Hazards and Risk Sources}
\paragraph{There is no single technique that works best in every situation. The most effective approach typically combines several methods, drawing on different sources of information and different modes of thinking. What follows is an overview of the most commonly used and most practically valuable techniques for compliance and legal professionals.}

\subsection{Structured Workshops and Brainstorming}
\paragraph{One of the most widely used approaches to risk identification is the \textbf{structured workshop} — bringing together a group of relevant stakeholders to identify risks collaboratively. Done well, a structured workshop combines the benefits of diverse perspectives, real-time discussion, and collective knowledge in a way that no desk-based analysis can replicate.}
\paragraph{The key word, however, is \textbf{\textit{structured}}. An unstructured brainstorming session — where participants simply list whatever comes to mind — tends to produce a list of well-known risks and recent concerns, whilst missing the less obvious or less comfortable ones. Effective workshop facilitation uses prompts, frameworks, and techniques to push participants beyond their default thinking.}
\paragraph{Useful facilitation approaches include:}
\begin{itemize}
    \bitem{Prompt questions}{\textit{"What keeps you awake at night?", "What would you be most embarrassed to explain to a regulator?", "What would happen if this key person left tomorrow?", "What assumption are we making that, if wrong, would cause serious harm?"}.}
    \bitem{Pre-mortems}{asking participants to imagine that a failure has already occurred and to work backwards to identify what caused it. This technique, borrowed from behavioural psychology, is particularly effective at surfacing risks that participants might otherwise be reluctant to raise.}
    \bitem{Scenario prompting}{presenting specific scenarios or categories of event and asking participants to consider whether and how they could affect the organisation}
    \bitem{Silent generation}{asking participants to write down risks individually before sharing them with the group, reducing the tendency for group thinking to suppress less popular or more uncomfortable contributions.}
\end{itemize}
\paragraph{For compliance professionals facilitating risk identification workshops, it is worth being deliberate about the composition of the group. Including people from different levels of the organisation — not just senior managers — and creating explicit space for quieter voices to be heard will produce a richer and more complete output.}

\subsection{Process Mapping and Workflow Analysis}
\paragraph{\textbf{Process mapping} involves documenting the sequence of steps involved in a particular activity or operation — and then systematically examining each step for potential failure points, vulnerabilities, and risk sources.}
\paragraph{This is a particularly valuable technique for identifying operational and compliance risks, because it grounds the risk identification in the reality of how work is actually done rather than how it is supposed to be done. The gap between the two is often a fertile source of risk — informal workarounds, shortcuts, and undocumented exceptions frequently represent genuine exposures that would not appear in a review of policy documentation alone.}
\paragraph{For each step in a process, useful questions include:}
\begin{itemize}
    \item{What could go wrong at this point?}
    \item{What happens if the person responsible makes a mistake, is absent, or lacks the necessary knowledge?}
    \item{What happens if a system or tool fails at this point?.}
    \item{What regulatory or legal requirements apply at this step, and how are they currently being met?}
    \item{What would a regulator or auditor look at first if they reviewed this step?.}
    \item{Are there any known near-misses or incidents associated with this part of the process?.}
\end{itemize}
\paragraph{Process mapping works particularly well when it involves the people who actually carry out the process — not just those who designed or manage it. The operational insight of frontline staff is often the most valuable input into this kind of analysis}

\subsection{Document and Data Review}
\paragraph{A significant amount of information relevant to risk identification already exists within most organisations — in incident logs, audit reports, regulatory correspondence, near-miss records, complaints data, staff surveys, and management information. Reviewing this material systematically is one of the most efficient and evidence-based approaches to risk identification.}
\paragraph{Key sources to consider include:}
\begin{itemize}
    \bitem{Incident and near-miss records}{past events are one of the most reliable indicators of future risk. Patterns in incident data can reveal systemic vulnerabilities that individual events might obscure.}
    \bitem{Internal audit reports}{audit findings frequently identify control weaknesses and process gaps that represent genuine risk sources.}
    \bitem{Regulatory correspondence and examination findings}{feedback from regulators is a direct and authoritative source of information about where your organisation's risk management is perceived to be inadequate.}
    \bitem{Complaints and customer feedback}{customer complaints are often early indicators of conduct risk, product risk, or service delivery failures.}
    \bitem{Management information and key risk indicators}{data on operational performance, control effectiveness, and risk metrics can reveal emerging trends that warrant closer investigation.}
    \bitem{Staff surveys and exit interviews}{employee feedback, particularly from those leaving the organisation, can surface cultural and people risks that other sources miss.}
\end{itemize}
\paragraph{One important caveat: the absence of incidents in a particular area is not the same as the absence of risk. An organisation that has not experienced a data breach does not necessarily have a low cybersecurity risk — it may simply have been fortunate. Document and data review should inform risk identification, not limit it.}

\subsection{Regulatory and Legal Horizon Scanning}
\paragraph{For compliance and legal professionals, horizon scanning — the systematic monitoring of developments in the regulatory and legal environment — is an essential component of risk identification. Regulatory change, new enforcement priorities, emerging legal obligations, and shifts in supervisory expectations can all introduce new risks or alter the significance of existing ones.}
\paragraph{Effective horizon scanning involves:}
\begin{itemize}
    \item{Monitoring regulatory publications, consultations, and guidance from relevant regulators and supervisory bodies}
    \item{Tracking legislative developments — both domestic and, where relevant, international — that may affect the organisation's compliance obligations.}
    \item{Following enforcement actions and regulatory decisions affecting peer organisations — these are often the clearest signals of where regulatory attention is focused and what standards are being applied.}
    \item{Engaging with industry bodies, legal advisers, and professional networks to stay abreast of emerging issues and developing expectations.}
    \item{Reviewing relevant case law and tribunal decisions that may affect the legal landscape within which the organisation operates.}
\end{itemize}
\paragraph{Horizon scanning should not be a one-off activity conducted at the beginning of a risk assessment cycle. It should be an ongoing process — feeding new information into the risk assessment framework on a continuous basis and triggering reassessment when significant changes occur.}

\subsection{Expert Consultation and External Benchmarking}
\paragraph{Some risks — particularly those involving technical, scientific, or highly specialised subject matter — are best identified with the assistance of external experts who bring knowledge and experience that the organisation does not have internally. Cybersecurity specialists, environmental consultants, actuaries, legal advisers in specialist areas, and regulatory consultants can all contribute valuable insight that internal teams may lack.}
\paragraph{\textbf{External benchmarking} — looking at the risk profiles and risk management experiences of comparable organisations — is also a valuable source of risk identification input. Industry associations, regulatory publications, published enforcement actions, and academic research can all shed light on risks that are common across a sector but may not yet have materialised in your own organisation.}
\paragraph{The caveat here is that external benchmarking should supplement, not replace, the internal risk identification process. Every organisation is different, and risks that are prominent in the sector may not be the most significant ones for your specific organisation — whilst risks that are unique to your circumstances may not appear in any external benchmark at all.}

\subsection{Checklists and Risk Libraries}
\paragraph{\textbf{Checklists} — pre-compiled lists of common risk sources and hazards relevant to a particular sector, activity, or risk type — can be useful as a prompt to ensure that the risk identification process does not overlook well-known categories of risk.}
\paragraph{Many regulatory bodies publish guidance that effectively constitutes a risk checklist — identifying the types of risk that the regulator expects firms to consider and address. Industry bodies and professional associations similarly publish risk catalogues and risk taxonomies that can serve as a useful starting point.}
\paragraph{The important caveat about checklists is that they are a floor, not a ceiling. Ticking off every item on a standard checklist without genuine engagement with the organisation's specific context is unlikely to produce a complete or useful risk identification. Checklists catch the obvious risks. The most consequential risks — the ones most likely to cause serious harm — are often the ones that do not appear on any standard list.}

%---------------------------- SECTION ----------------------------
\section{Organising What You Find — From Identification To Risk List}
\paragraph{As the risk identification process generates outputs — from workshops, document reviews, process maps, and other sources — those outputs need to be organised into a structured form that can feed into the evaluation step that follows.}
\paragraph{At this stage, the goal is not to evaluate risks — it is simply to capture them clearly and completely. A useful approach is to record each identified risk in a consistent format that capture:}
\begin{itemize}
    \bitem{A clear description of the risk}{what could happen, and in what circumstances.}
    \bitem{The risk category}{mapping the risk to the taxonomy established at the outset, drawing on the categories explored in Chapter~\ref{chap:Types of Risk}.}
    \bitem{The risk source}{what drives this risk? Is it internal or external? Is it driven by people, processes, systems, or external events?.}
    \bitem{The potential consequences}{at this stage, a high-level description of what the impact could be if the risk materialised — financial, regulatory, reputational, operational, or human.}
    \bitem{Any existing controls}{a note of any measures already in place that are designed to manage this risk, which will be important in the evaluation step.}
\end{itemize}
\paragraph{This structured capture process begins the construction of the risk register — the central documentary artefact of the risk assessment process, which we will develop further in Chapter~\ref{chap:Step 3 — Evaluating and Analysing Risks} and Chapter~\ref{chap:Step 5 - Recording, Reviewing, and Monitoring}.}

%---------------------------- SECTION ----------------------------
\section{Common Pitfalls in Risk Identification}
\paragraph{Having explored the techniques for effective risk identification, it is worth being equally explicit about the pitfalls — the patterns of failure that recur consistently across organisations and sectors.}

\subsection{Availability Bias}
\paragraph{Humans — and organisations — tend to give disproportionate weight to risks that are recent, vivid, or memorable, and to underweight risks that are less salient even if they are equally or more significant. This \textbf{availability bias} means that risk identification exercises conducted shortly after a major incident tend to focus heavily on that type of incident, whilst risks that have not recently materialised may be systematically underweighted.}
\paragraph{The antidote is deliberate structure — using techniques and frameworks that ensure all categories of risk receive systematic attention, rather than allowing the process to be driven by whatever happens to be most prominent in people's minds at the time.}

\subsection{Groupthink}
\paragraph{In workshop settings, the natural social dynamics of groups can suppress dissenting views, unconventional thinking, and uncomfortable observations. When senior people in the room have already expressed a view, other participants may be reluctant to challenge it. When a consensus emerges, the natural tendency is to move on rather than probe further.}
\paragraph{Effective facilitation — including anonymous input mechanisms, pre-mortem techniques, and explicit encouragement of challenge — can mitigate groupthink. So can ensuring that risk identification workshops include a genuinely diverse group of participants, rather than a homogeneous team of senior managers.}

\subsection{Optimism Bias}
\paragraph{There is a well-documented human tendency to underestimate the likelihood of negative events and to overestimate our ability to prevent or manage them. In an organisational context, this \textbf{optimism bias} can manifest as a systematic tendency to rate risks as less likely or less serious than they actually are — particularly when those risks relate to the organisation's own processes, people, or decisions.}
\paragraph{Encouraging participants to think about risks from the perspective of a regulator, an auditor, or a sceptical external observer — rather than from the perspective of someone who believes in the organisation's ability to manage things well — can help counteract this tendency.}

\subsection{The Silence of the Familiar}
\paragraph{Some of the most significant risks in any organisation are so familiar, so embedded in normal operations, that they have become invisible. The process that everyone knows is unreliable but has always worked so far. The system that is well past its replacement date but has never actually failed. The key person whose departure would be catastrophic but whose knowledge is not documented anywhere.}
\paragraph{These risks do not surface easily in workshops or document reviews precisely because they are taken for granted. Identifying them requires deliberate effort — specifically looking for the things that are known to be fragile, the dependencies that are understood but not acknowledged, and the assumptions that underpin normal operations without ever being tested.}

\subsection{Scope Fixation}
\paragraph{As noted in Chapter~\ref{chap:The Risk Assessment Process — An Overview}, scope is a necessary constraint on any risk assessment. But there is a risk that the scope boundary — particularly if drawn early and without sufficient thought — becomes a mental barrier that prevents participants from identifying risks that sit just outside it. Third-party risks, systemic risks, and risks that span organisational boundaries are particularly susceptible to being missed for this reason.}
\paragraph{Building in a deliberate prompt to consider risks at and beyond the scope boundary — and being willing to adjust the scope if genuinely significant risks are identified outside it — is good practice.}

%---------------------------- SECTION ----------------------------
\section{Human Factors and Organisational Blind Spots}
\paragraph{No review of risk identification would be complete without acknowledging the role of human factors — the ways in which the behaviour, cognition, and culture of the people within an organisation shape its risk profile.}
\paragraph{Human factors sit behind a significant proportion of the most serious organisational failures. Mistakes, misunderstandings, poor judgement, wilful non-compliance, cultural pressures, fatigue, and inadequate training all contribute to risk in ways that are often under-represented in formal risk identification processes that focus primarily on systems, processes, and external events.}
\paragraph{For compliance professionals, several human factor risks deserve particular attention:}
\begin{itemize}
    \bitem{Conduct risk}{the risk that individuals will behave in ways that harm customers, markets, or the organisation, whether deliberately or through poor judgement.}
    \bitem{Competence risk}{the risk that individuals lack the knowledge, skills, or experience to perform their responsibilities to the required standard.}
    \bitem{Cultural risk}{the risk that the organisation's values, norms, and incentive structures encourage or permit behaviour that creates legal, regulatory, or ethical exposure.}
    \bitem{Change fatigue}{in organisations that have experienced significant change, the risk that people are overwhelmed, disengaged, or cutting corners increases substantially.}
    \bitem{Incentive misalignment}{where individual incentives are not aligned with the organisation's risk management objectives, people may take on more risk than the organisation intends, or may be motivated to conceal problems rather than surface them.}
\end{itemize}
\paragraph{Identifying human factor risks requires a degree of honesty and self-awareness that can be uncomfortable. It means being willing to ask hard questions about culture, leadership behaviour, incentive structures, and the gap between stated values and actual practice. But it is precisely because these questions are uncomfortable that they are so important — and so valuable — to ask.}

%---------------------------- SECTION ----------------------------
%\newpage
\section{Chapter Summary}
In this chapter we learned about the following key points:
\begin{table}[H]
  \centering
  \renewcommand{\arraystretch}{1.5}
  \begin{tabularx}{\textwidth}{ | >{\raggedright\arraybackslash}p{0.35\textwidth} 
                                | >{\raggedright\arraybackslash}X | }
    \hline
    \textbf{Technique} & \textbf{Best Used For} \\
    \hline
    \textbf{Structured workshops} & Drawing on diverse perspectives; surfacing subjective and cultural risks\\
    \hline
    \textbf{Process mapping} & Identifying operational, process, and compliance risks in specific activities\\
    \hline
    \textbf{Document and data review} & Evidence-based identification drawing on incident history and audit findings\\
    \hline
    \textbf{Horizon scanning} & Identifying regulatory, legal, and environmental changes that introduce new risk\\
    \hline
    \textbf{Expert consultation} & Technical, specialist, or highly complex risk areas\\
    \hline
    \textbf{External benchmarking} & Understanding sector-wide risk patterns and regulatory expectations\\
    \hline
    \textbf{Checklists and risk libraries} & Ensuring completeness — as a prompt, not a substitute for genuine analysis\\
    \hline
  \end{tabularx}
  \caption{Identification Techniques}
\end{table}

\begin{table}[H]
  \centering
  \renewcommand{\arraystretch}{1.5}
  \begin{tabularx}{\textwidth}{ | >{\raggedright\arraybackslash}p{0.35\textwidth} 
                                | >{\raggedright\arraybackslash}X | }
    \hline
    \textbf{Pitfall} & \textbf{What It Looks Like} \\
    \hline
    \textbf{Availability bias} & Over-focusing on recent or vivid risks at the expense of less salient ones\\
    \hline
    \textbf{Groupthink} & Social dynamics suppressing unconventional or uncomfortable observations\\
    \hline
    \textbf{Optimism bias} & Systematically underestimating likelihood or severity of negative events\\
    \hline
    \textbf{Silence of the familiar} & Known fragilities being overlooked because they have always worked so far\\
    \hline
    \textbf{Scope fixation} & Failing to identify risks that sit at or beyond the defined scope boundary\\
    \hline
  \end{tabularx}
  \caption{Pitfalls}
\end{table}

\paragraph{Key takeaways:}
\begin{itemize}
  \item{Risk identification is the most consequential step in the process — everything that follows depends on its quality.}
  \item{No single technique is sufficient — the most effective approach combines several methods drawing on different sources and perspectives.}
  \item{The conditions for good risk identification — psychological safety, diverse perspectives, and structured curiosity — are as important as the techniques themselves.}
  \item{Human factors and organisational blind spots are among the most significant and most consistently underidentified categories of risk.}
  \item{The output of this step is a structured risk list — the foundation of the risk register that will be developed through the steps that follow.}
  \item{Checklists and frameworks are a floor, not a ceiling — the most important risks are often the ones that do not appear on any standard list.}
\end{itemize}

\barquote{In Chapter~\ref{chap:Step 2 — Determining Who or What Is Affected}, we move to the second step of the process — determining who or what is affected by the risks we have identified, and why thinking carefully about exposure is more important, and more nuanced, than it might initially appear.}